\newpage
\section{Validity of Human Agreement Guarantee}
\subsection{Proof of Theorem \ref{thm:1}} \label{app:proof_of_theorem1}
\begin{customthm}{1}
    Consider a threshold $\widehat{\lambda}$ chosen as in \S \ref{sec:providing_human_agreement_guarantee}, and a selective evaluator $(f_\textit{LM}, c_\textit{LM})$ operating based on $\widehat{\lambda}$. Then, Equation (\ref{eq:human-agreement-guarantee}) is satisfied with probability at least $1 - \delta$.
\end{customthm}

The proof extends that of Theorem B.1 in \cite{rcps}. Let $R(\lambda)$ denote the true risk of disagreeing with humans at threshold $\lambda$. It suffices to show that $P(R(\widehat{\lambda}) \leq \alpha) \geq 1 - \delta$. We first note that $n(\lambda)\widehat{R}(\lambda)$ is a binomial random variable, \ie
\begin{equation*}
    n(\lambda)\widehat{R}(\lambda) \sim \bin(n(\lambda), R(\lambda)).
\end{equation*}
Thus, the lower tail bound for $\widehat{R}(\lambda)$ can be expressed as a function $g$ of $t \in \mathbb{R}$ and $R(\lambda)$ as
\begin{equation*}
    P(\widehat{R}(\lambda) \leq t) = P\big(\bin(n(\lambda), R(\lambda)) \leq \lceil n(\lambda)t \rceil \big) \eqqcolon g(t; R(\lambda)).
\end{equation*}
Plugging this into the definition of $\widehat{R}^+(\lambda)$ in Equation \ref{eq:r-hat},
\begin{align*}
    \widehat{R}^+(\lambda) & = \sup \left\{ R(\lambda): P\big(\bin(n(\lambda), R(\lambda)) \leq \lceil n(\lambda) \widehat{R}(\lambda) \rceil \big) \geq \delta \right\} \\
    & = \sup \left\{ R(\lambda): g(\widehat{R}(\lambda); R(\lambda)) \geq \delta \right\}.
\end{align*}
Here, let G denote the CDF of $\widehat{R}(\lambda)$ and $G^{-1}(\delta) = \sup \{ x: G(x) \leq \delta \}$. From above, we know that if $R(\lambda) > \widehat{R}^+(\lambda)$, then $g(\widehat{R}(\lambda); R(\lambda)) < \delta$. Therefore,
\begin{align*}
    P(R(\lambda) > \widehat{R}^+(\lambda)) & \leq P(g(\widehat{R}(\lambda); R(\lambda)) < \delta) \\
    & = P(G(\widehat{R}(\lambda)) < \delta) \\
    & \leq P(\widehat{R}(\lambda) < G^{-1}(\delta)) \\
    & \leq \delta.
\end{align*}
Hence, $P(R(\lambda) \leq \widehat{R}^+(\lambda)) \geq 1 - \delta$, implying that $\widehat{R}^+(\lambda)$ is the $(1 - \delta)$ upper confidence bound of $R(\lambda)$. Finally, since we have $\widehat{R}^+(\widehat{\lambda}) \leq \alpha$ from the definition of $\widehat{\lambda}$, we obtain
\begin{equation*}
    P(R(\widehat{\lambda}) \leq \widehat{R}^+(\widehat{\lambda}) \leq \alpha) \geq 1 - \delta. \qed
\end{equation*}

\subsection{Extension to Cascades of Judge Models} \label{app:cascades_of_judge_models}
\begin{algorithm}[b]
\caption{$\,\calibrate(\mathcal{M}, D_\textit{cal}, \alpha, \delta)$}\label{alg:calibrate}
\begin{algorithmic}
\Require A list of judges $\mathcal{M} = (M_1, \cdots, M_{|\mathcal{M}|})$, a calibration set $D_\textit{cal}$ and test set $D_\textit{test}$ to be evaluated, risk tolerance $\alpha$ and error level $\delta$
\Ensure A set of calibrated thresholds $\Lambda$
\State $\Lambda \gets \emptyset$ \Comment{Initialize the set of thresholds.}
\State $D \gets D_\textit{cal}$ \Comment{Initialize $D$ for calibrating each model.}
\For{$i = 1 \text{ to } |\mathcal{M}|$}
    \State $\lambda_i \gets \calibratesingle(M_i, D, \alpha, \frac{\delta}{|\mathcal{M}|})$ \Comment{Calibrate $\lambda_i$ for each model $M_i$.}
    \State $\Lambda \gets \Lambda \cup \{\lambda_i\}$
    \State $D \gets \{(x, y) \in D: c_{M_i}(x) < \lambda_i\}$ \Comment{Update $D$ with only the previously abstained instances.}
\EndFor
\Return $\Lambda$ \Comment{Return the set of calibrated thresholds.}
\end{algorithmic}
\end{algorithm}

We illustrate the extension of our test procedure from a single model to the cascades of judge models. For notational simplicity, we denote the test procedure for a single judge model in \S \ref{sec:providing_human_agreement_guarantee} as a function $\calibratesingle$. This function takes as input a model $M$, a calibration set $D$, risk tolerance $\alpha$ and error level $\delta$, and gives a calibrated threshold $\widehat{\lambda}$ as an output.

The calibration procedure for cascades of judge models is shown in Algorithm \ref{alg:calibrate}. The procedure sequentially applies $\calibratesingle$ to each judge model. Specifically for each judge model $M_i$, we calibrate $\lambda_i$ by testing over the set of instances $D$ that have been abstained by the previous models. This allows us to ensure that for each $M_i$,
\begin{equation*}
    P\left(f_{M_i}(x) = y_\textit{human} \Bigg| c_{M_i}(x) \ge \lambda_i, \bigwedge_{j=1}^{i-1} c_{M_j}(x) < \lambda_j\right) \ge 1 - \alpha
\end{equation*}
is satisfied with probability at least $1 - \frac{\delta}{|\mathcal{M}|}$. To provide the guarantee across all judge models, define $R_i$ as the disagreement risk for each judge model $M_i$:
\begin{equation*}
    R_i \coloneqq P\left(f_{M_i}(x) \neq y_\textit{human} \Bigg| c_{M_i}(x) \ge \lambda_i, \bigwedge_{j=1}^{i-1} c_{M_j}(x) < \lambda_j\right).
\end{equation*}
Also, define $R_\textit{cascades}$ as the risk of full cascaded selective evaluation, \ie
\begin{equation*}
    R_\textit{cascades} \coloneqq P(f_\textit{cascades}(x) \neq y_\textit{human} |\,x \textit{ not abstained by the cascades}).
\end{equation*}
It is easy to see that $R_\textit{cascades}$ is an interpolation between all $R_i$s, thus $R_\textit{cascades} \leq \max_i R_i$. Therefore,
\begin{align*}
    P\big(R_\textit{cascades} > \alpha \big) & \leq P\big(\max_i R_i > \alpha\big) = P\bigg(\bigcup_i R_i > \alpha \bigg) \leq \sum_i P\big(R_i > \alpha \big),
\end{align*}
where the last inequality comes from union bound. Since we know that $P(R_i > \alpha) \leq \frac{\delta}{|\mathcal{M}|}$ for each judge model $M_i$, 
\begin{equation*}
    \sum_i P\big(R_i > \alpha \big) \leq \frac{\delta}{|\mathcal{M}|} \cdot |\mathcal{M}| = \delta.
\end{equation*}
Thus, $P(R_\textit{cascades} > \alpha) \leq \delta$. In other words, the risk of disagreement across all judge models is guaranteed to be at most $\alpha$, with probability at least $1 - \delta$.

 % \\
 %    & \leq \frac{\delta}{|\mathcal{M}|} \\
 %    & = \delta.

\newpage
\section{Additional Results on Confidence Estimation}\label{app:additional_results_on_confidence_estimation}
\begin{table*}[h]\centering
\vspace{-5pt}
\caption{Additional results on confidence estimation with Mistral-7B. We find that more sophisticated methods that measure the semantic variance between chain-of-thoughts often underperform Simulated Annotators, marking similar or worse performance with zero-shot predicted probability.}
\resizebox{.95\textwidth}{!}{
    \begin{tabular}{clcccccccc}\toprule
        \multicolumn{2}{c}{\textbf{Dataset}} & \multicolumn{4}{c}{\textbf{AlpacaEval}}& \multicolumn{4}{c}{\textbf{TL;DR}} \\\cmidrule(lr){3-6}\cmidrule(lr){7-10}
        \multicolumn{2}{c}{\textbf{Method}} & Acc. & ECE $\downarrow$ & AUROC & AUPRC & Acc. & ECE $\downarrow$ & AUROC & AUPRC \\
        \midrule
        \multirow{5}{*}[-0.3em]{\shortstack{\textit{Mistral-}\\\textit{7B-it}}} & Predictive Probability (CoT) & 0.636 & 0.292 & 0.527 & 0.655 & 0.669 & 0.275 & 0.637 & 0.751 \\
                                      & Lexical Similarity & 0.636 & 0.478 & 0.478 & 0.623 & 0.669 & 0.459 & 0.520 & 0.639 \\
                                      & Semantic Sets & 0.636 & - & 0.513 & 0.638 & 0.669 & - & 0.545 & 0.665 \\
                                      & Semantic Entropy & 0.636 & - & 0.572 & 0.684 & 0.669 & - & 0.650 & 0.762 \\
        \cmidrule(r){2-10}
                                      & Simulated Annotators (Ind.) & \textbf{0.684} & \textbf{0.075} & \textbf{0.632} & \textbf{0.772} & \textbf{0.696} & \textbf{0.103} & \textbf{0.654} & \textbf{0.807} \\
        \bottomrule
    \end{tabular}
}
\label{tab:additional_confidence_estimation}
\end{table*}

In Table \ref{tab:additional_confidence_estimation}, we provide additional results for more sophisticated methods for confidence estimation. Given an instance $x$, these methods first generate $M$ chain-of-thoughts (CoTs) from an LLM judge prior to inferring the preference label, then measure their variance either in the label space or on the semantic-level:
\begin{itemize}
    \item \textit{Predictive Probability (CoT)}: We average the $M$ label predictive probabilities assigned by the LLM judge after generating chain-of-thoughts.
    \item \textit{Lexical Similarity}: As a simple proxy of semantic variance, we average ROUGE-L across all pairs of $M$ chain-of-thoughts. The intuition is that when the CoTs exhibit high lexical overlap with each other, the model is relatively confident about its generation.
    \item \textit{Semantic Sets} \citep{semantic-eigenvalue}: We cluster the CoTs into semantically equivalent groups using a bidirectional entailment model \citep{alignscore}, then use the number of clustered groups to represent model uncertainty.
    \item \textit{Semantic Entropy} \citep{semantic-entropy}: We additionally use the likelihood of each generated chain of thought, and measure the average sequence-level entropy across the semantically equivalent groups.
\end{itemize}
We follow \cite{semantic-eigenvalue} to set $M = 20$. For \textit{Semantics Sets} and \textit{Semantic Entropy}, we exclude expected calibration error as the two scores represent model uncertainty rather than the confidence score calibrated in $[0, 1]$. These methods incur significant overhead compared to the methods discussed in \S \ref{sec:simulated_annotators}, generating 20 sequences of chain-of-thoughts for each input instance and employing a supervised entailment model. Nonetheless, we find that their performance consistently underperforms that of \textit{Simulated Annotators}, with the best method \textit{Predictive Probability (CoT)} performing worse than our ablation \textit{Randomized Annotators} (Table \ref{tab:simulated_annotators_main}).

\newpage
\section{Additional Results on Chat Assistant Evaluation} \label{app:additional_results_on_chat_assistant_evaluation}
\begin{table*}[h]\centering
\caption{Comparison to baselines on Auto-J, with target agreement level $1 - \alpha = 0.8$. The results are averaged across 1000 runs with random data split.}
\resizebox{.96\textwidth}{!}{
    \begin{tabular}{lcccccc}\toprule
         \multirow{2}{*}[-0.3em]{\textbf{Method}} & \multicolumn{3}{c}{\textbf{Evaluator Composition (\%)}} & \multirow{2}{*}[-0.3em]{\textbf{Coverage (\%)}} & \multirow{2}{*}[-0.3em]{\shortstack{\textbf{Guarantee Success}\\\textbf{Rate (\%)}}}\\\cmidrule(lr){2-4}
         & \textit{Mistral-7B} & \textit{GPT-3.5-turbo} & \textit{GPT-4-turbo} & & \\
         \midrule
         No Selection & 0.0 & 0.0 & 100.0 & 100.0 & 0.0 \\
         Heuristic Selection & 0.0 & 0.0 & 100.0 & 72.0 & 2.4 \\
         Cascaded Heuristic Selection & 36.6 & 40.2 & 23.2 & 84.4 & 0.8 \\
         \midrule
         \multirow{3}{*}{Point-Estimate Calibration} & 100.0 & 0.0 & 0.0 & 12.1 & 54.3 \\
                                                     & 0.0 & 100.0 & 0.0 & 27.9 & 52.1 \\
                                                     & 0.0 & 0.0 & 100.0 & 42.9 & 56.5 \\
         \midrule
         \textbf{Cascaded Selective Evaluation} & \textbf{49.3} & \textbf{21.8} & \textbf{28.9} & \textbf{42.6} & \textbf{90.2} \\
         \bottomrule
    \end{tabular}
}
% \vspace{-5pt}
\label{tab:auto-j_results}
\end{table*}
\begin{figure*}[h]
    \centering
    \includegraphics[width=.96\textwidth]{figures/auto-j_results.pdf}
    \vspace{-5pt}
    \caption{Human Agreement Guarantee on Auto-J. \textbf{GPT-4 without abstention obtains only 63.2\% agreement with humans, while \method guarantees target human agreement level of up to 80\% with high probability.}}
    \vspace{-10pt}
    \label{fig:auto-j_results}
\end{figure*}

\section{Evaluation under Distribution Shift} \label{app:evaluation_under_distribution_shift}
\begin{table*}[h]\centering
\caption{Evaluation under distribution shift on ChatArena. We induce distribution shift by sampling the calibration and test set respectively from two disjoint sets of instances with no overlap of evaluated models. We iterate experiments for 1000 random splits and aggregate the results. \textbf{\method guarantees high agreement even under the realistic distribution shift.}}
\resizebox{.75\textwidth}{!}{
    \begin{tabular}{cccc}\toprule
         \multirow{2}{*}{\shortstack{\textbf{ Target Human }\\\textbf{ Agreement (\%) }}} & \multirow{2}{*}{\shortstack{\textbf{ Empirical Human }\\\textbf{ Agreement (\%) }}} & \multirow{2}{*}{\textbf{ Coverage (\%) }} & \multirow{2}{*}{\shortstack{\textbf{ Guarantee Success }\\\textbf{ Rate (\%) }}} \\
         & & & \\
         \midrule
         70.0 & 73.4 & 100.0 & 100.0 \\
         75.0 & 75.3 & 91.4 & 92.5 \\
         80.0 & 80.8 & 72.1 & 90.8 \\
         85.0 & 85.2 & 55.4 & 91.0 \\
         90.0 & 90.1 & 31.8 & 90.7 \\
         \bottomrule
    \end{tabular}
}
\label{tab:evaluation_under_distribution_shift}
\end{table*}

\newpage
\section{Details on Experimental Setup} 
\subsection{Prompts}
\begin{tcolorbox}[colback=white,colframe=black!75!white,colbacktitle=black!75!white,title=Prompt for Summarization Evaluation]
  You are a helpful assistant.\\

  Given a document and two summaries of the document, an annotator chose which summary is preferred. Given examples of the annotator's decision, predict the annotator's verdict on the given example. If Summary A is preferred to Summary B, the annotator chose ``[[A]]''. If Summary B is preferred to Summary A, the annotator chose ``[[B]]''.\\

  [Document]\\
  \{document\_\_example\_1\}\\

  [Summary A]\\
  \{summary\_a\_example\_1\}\\

  [Summary B]\\
  \{summary\_b\_example\_1\}\\

  [Verdict]:\\
  \{verdict\_example\_1\}\\

  \texttt{...[few-shot examples omitted]}\\

  [Document]\\
  \{document\}\\

  [Summary A]\\
  \{summary\_a\}\\

  [Summary B]\\
  \{summary\_b\}\\

  Verdict (either ``[[A]]'' or ``[[B]]''):
\end{tcolorbox}
\newpage
\begin{tcolorbox}[colback=white,colframe=black!75!white,colbacktitle=black!75!white,title=Prompt for Chat Assistant Evaluation]
  You are a helpful assistant.\\

  Given a question and two assistant's answers to the question, an annotator chose which assistant's answer is preferred. Given examples of the annotator's decision, predict the annotator's verdict on the given example. If Assistant A's response is preferred to Assistant B's, the annotator chose ``[[A]]''. If Assistant B's response is preferred to Assistant A's, the annotator chose ``[[B]]''.\\

  [Question]\\
  \{question\_example\_1\}\\

  [Assistant A's response]\\
  \{assistant\_a\_example\_1\}\\

  [Assistant B's response]\\
  \{assistant\_b\_example\_1\}\\

  [Verdict]:\\
  \{verdict\_\_example\_1\}\\

  \texttt{...[few-shot examples omitted]}\\

  [Question]\\
  \{question\}\\

  [Assistant A's response]\\
  \{assistant\_a\}\\

  [Assistant B's response]\\
  \{assistant\_b\}\\

  Verdict (either ``[[A]]'' or ``[[B]]''):
\end{tcolorbox}

\newpage
\subsection{Human Evaluation Details} \label{app:human_evaluation_details}
\begin{figure*}[h]
    \centering
    \includegraphics[width=1.0\linewidth]{figures/annotation_ui.png}
    \caption{Human Annotation Interface.}
    \label{fig:annotation_ui}
\end{figure*}
\paragraph{Annotator Recruitment.} We recruited annotators from Prolific\footnote{\url{https://app.prolific.com}} who have recorded at least 99\% approval rate, are fluent in English, and have completed a Bachelor's degree. In addition, we manually designed 10 qualification examples based on our annotation guidelines. The purpose of the qualification test is to find annotators who understand and carefully follow our guidelines. Participants who scored more than 80\% were included in our actual human study. We qualified 21 annotators to do the study and paid them \$15 per hour. 

\paragraph{Annotation Task.} We randomly sample 600 examples from ChatArena~\citep{mt-bench} each consisting of a query and two model responses. Given each example, we instruct the annotators to select the overall better response considering several aspects such as helpfulness, truthfulness, and harmlessness. Each instance is evaluated by 3-5 annotators. We also allow annotators to occasionally skip an instance with a reason if they have no idea how to evaluate it. We provide the screenshot of our annotation guidelines and interface in Figure \ref{fig:annotation_ui} and \ref{fig:annotation-guidelines}.



\newpage
\begin{figure*}[h]
    \centering
    \includegraphics[width=0.9\linewidth]{figures/annotation_guidelines.pdf}
    \caption{Human Annotation Guidelines.}
    \label{fig:annotation-guidelines}
\end{figure*}

\newpage
\section{Qualitative Examples}
\vspace{10pt}
\begin{tcolorbox}[colback=white,colframe=black!75!white,colbacktitle=black!75!white,title=Summarization Example 1]
  \textbf{Document} \\
  Long story short -- We dated for about a year and 1/2 I moved for my job we couldn't handle the distance and we end our relationship. Months go by, I see him Friday night and he is very polite and nice and after we started drinking he started pouring his heart out to me. The next day he calls me twice in the morning and I missed the calls, try to call him back. Doesnt answer, wait a few days, send him a long message, and it said he read it, and has still said nothing. Should I pretend like it never happen?
  \tcbline
  \textbf{Summary A} \\
  Ex-bf of 1.5 years went full blown confessional on me this weekend. Should I pretend like it never happened and pretend like nothing ever happened?
  \tcbline
  \textbf{Summary B} \\
  Ex-bf of 1.5 years confided in me. Called me twice in the morning and I missed the calls. Doesnt answer my messages. Should I pretend to ignore it?
  \tcbline
  \textbf{Evaluator}: \textit{Abstained} \quad \textbf{Model Preference}: \textit{None} \quad \textbf{Human Preference}: \textit{B is better}\\
  \textbf{Inter-Annotator Agreement}: $0.75$
\end{tcolorbox}
\vspace{10pt}
\begin{tcolorbox}[colback=white,colframe=black!75!white,colbacktitle=black!75!white,title=Summarization Example 2]
  \textbf{Document} \\
  I have been working since 16 and make around \$12,000 annually with a small amount of expenses (gas, car insurance). I literally just turned 18 and felt that it would be a good idea to get a credit card. I honestly don't need it to ``make purchases I can't afford'' but I really just want one to start building my credit.
  
  I talked with my Dad about it and he said that applying for a store card like Best Buy (where I work actually) would be a good idea or that getting a secured credit card would also be a good idea. I prefer the secured credit card route due to being able to use it anywhere and I can't see myself actually using a store card consistently as I do not shop at one place consistently.
  
  What are some of the best secured Credit Cards to get? I'm not worried about a high credit limit, I just want my credit to go up and be good, for a lack of better words. My current debit card is with a Credit Union and I read that those and Capital One are usually the best route to go with.
  
  I also wanted to know if it is better to with Visa, Mastercard, Discover, AMEX ETC.
  \tcbline
  \textbf{Summary A} \\
  18 y/o male wants to start building credit and wants a secure credit card to do so. Willing to use anywhere, not worried about credit limit. Best secured credit cards?
  \tcbline
  \textbf{Summary B} \\
  just turned 18 wanting to build credit. Talked to Dad about applying for secured credit card or Best Buy store card. Best route?
  \tcbline
  \textbf{Evaluator}: \textit{Mistral-7B} \quad \textbf{Model Preference}: \textit{A is better} \quad \textbf{Human Preference}: \textit{A is better}\\
  \textbf{Inter-Annotator Agreement}: $0.6$
\end{tcolorbox}

\newpage
\begin{tcolorbox}[colback=white,colframe=black!75!white,colbacktitle=black!75!white,title=Chat Assistant Example 1]
  \textbf{Query} \\
  Who is Nicola Tesla?
  \tcbline
  \textbf{Assistant A} \\
  A great electrical engineer
  \tcbline
  \textbf{Assistant B} \\
  Nicola Tesla was an electrical engineer, a physicist, and a futurist. She was an early discoverer of electromagnetic radiation (EM) and the first person to use radio waves to transmit information. Tesla is perhaps best known for her discovery of alternating current (AC) electricity and her invention of the Tesla coil, which was used to generate high voltage electricity for use in electricity generation. However, Tesla is also a controversial figure, with some historians and scientists questioning her work and reputation.
  \tcbline
  \textbf{Evaluator}: \textit{Mistral-7B} \quad \textbf{Model Preference}: \textit{B is better} \quad \textbf{Human Preference}: \textit{B is better}\\
  \textbf{Inter-Annotator Agreement}: $1.0$
\end{tcolorbox}
\vspace{10pt}
\begin{tcolorbox}[colback=white,colframe=black!75!white,colbacktitle=black!75!white,title=Chat Assistant Example 2]
  \textbf{Query} \\
   If ``aka blux fii'' means ``the cat is white'', ``aka jorr fii'' means ``the chicken is white'', ``aka urf fii'' means ``the milk is white'' and ``urf blux orto'' means ``the cat likes milk'', how would you say ``the cat likes chicken''?
  \tcbline
  \textbf{Assistant A} \\
  If ``aka blux fii'' means ``the cat is white'', ``aka jorr fii'' means ``the chicken is white'', ``aka urf fii'' means ``the milk is white'', and ``urf blux orto'' means ``the cat likes milk'', then to say ``the cat likes chicken'', you would say ``aka jorr fii.''
  \tcbline
  \textbf{Assistant B} \\
  To say ``the cat likes chicken,'' you would use ``blux jorr orto.'' Here's the breakdown:
  
  - ``aka'' means ``the'' \\
  - ``blux'' means ``cat'' \\
  - ``jorr'' means ``chicken'' \\
  - ``fii'' means ``is white'' \\
  - ``urf'' means ``milk'' \\
  - ``orto'' means ``likes'' \\
  
  So, ``blux jorr orto'' translates to ``the cat likes chicken''.
  \tcbline
  \textbf{Evaluator}: \textit{GPT-4} \quad \textbf{Model Preference}: \textit{B is better} \quad \textbf{Human Preference}: \textit{B is better}\\
  \textbf{Inter-Annotator Agreement}: $1.0$
\end{tcolorbox}

\newpage
\begin{tcolorbox}[colback=white,colframe=black!75!white,colbacktitle=black!75!white,title=Chat Assistant Example 3]
  \textbf{Query} \\
  Privacy is not a universal right.
  \tcbline
  \textbf{Assistant A} \\
  Privacy is not a universal right. From an international human rights perspective, privacy is not a human right because it is not considered a fundamental right. However, the American Civil Liberties Union (ACLU) has highlighted that in the US, the right to privacy is considered a fundamental right under the US Constitution's Fourth Amendment.
  \tcbline
  \textbf{Assistant B} \\
  That statement is not entirely accurate. While some countries may not have a specific constitutional guarantee of privacy, the right to privacy is often recognized as a fundamental human right under international law and is protected by various treaties and conventions, such as the Universal Declaration of Human Rights and the European Convention on Human Rights. Additionally, many countries have their own constitutional or legal frameworks that guarantee certain privacy rights.
  \tcbline
  \textbf{Evaluator}: \textit{GPT-3.5} \quad \textbf{Model Preference}: \textit{B is better} \quad \textbf{Human Preference}: \textit{B is better}\\
  \textbf{Inter-Annotator Agreement}: $1.0$
\end{tcolorbox}
\vspace{10pt}
\begin{tcolorbox}[colback=white,colframe=black!75!white,colbacktitle=black!75!white,title=Chat Assistant Example 4]
  \textbf{Query} \\
  list all yearly hottest temperature records starting from 1900 to the present
  \tcbline
  \textbf{Assistant A} \\
  Here is a list of the hottest temperatures recorded each year starting from 1900 to the present:
  \begin{itemize}[leftmargin=*]
      \item 1900 - Death Valley, California at 134\textdegree F (56.7\textdegree C) on July 10, 1913
      \item 1901 -  Aziziyah, Libya at 136\textdegree F (58\textdegree C) on September 13, 1922
      \item 1903 - Death Valley, California at 134\textdegree F (56.7\textdegree C) on July 10, 1913
  \end{itemize}
  ...\texttt{(omitted)}
  
  [The list continues up to 2020]
  \tcbline
  \textbf{Assistant B} \\
  I'm sorry, but as an AI language model, I don't have access to real-time data. However, I can provide you with a list of some of the most notable hottest temperature records starting from 1900 to the present, based on historical records and news reports:\\[3pt]
  1. 1936 - 134\textdegree F (57\textdegree C) in Death Valley, California \\[3pt]
  2. 1913 - 134\textdegree F (57\textdegree C) in Turkey \\[3pt]
  3. 2020 - 135\textdegree F (62.8\textdegree C) in Oman \\[3pt]
  ...\texttt{(omitted)}\\[3pt]
  10. 2021 - 120.6\textdegree F (49.2\textdegree C) in Kuwait City, Kuwait
  \tcbline
  \textbf{Evaluator}: \textit{Abstained} \quad \textbf{Model Preference}: \textit{None} \quad \textbf{Human Preference}: \textit{A is better} \\
  \textbf{Inter-Annotator Agreement}: $0.75$
\end{tcolorbox}

\newpage
\begin{tcolorbox}[colback=white,colframe=black!75!white,colbacktitle=black!75!white,title=Chat Assistant Example 5]
  \textbf{Query} \\
  What month and year will the US stock market make a bottom?
  \tcbline
  \textbf{Assistant A} \\
  The US stock market has been in a bear market since February 2018. The DJIA and S\&P 500 have both made lows in that time. The S\&P 500 has made a low on March 1220 and the DJIA on February 24, 2019.
  \tcbline
  \textbf{Assistant B} \\
  As an AI, I am unable to predict future events, especially those related to the stock market. Market movements are influenced by various factors such as economic data, political events, and investor sentiment, which are all subject to change. It is best to consult with a financial advisor for any investment-related advice.
  \tcbline
  \textbf{Evaluator}: \textit{GPT-4} \quad \textbf{Model Preference}: \textit{B is better} \quad \textbf{Human Preference}: \textit{B is better}\\
  \textbf{Inter-Annotator Agreement}: $0.83$
\end{tcolorbox}
\vspace{10pt}
\begin{tcolorbox}[colback=white,colframe=black!75!white,colbacktitle=black!75!white,title=Chat Assistant Example 6]
  \textbf{Query} \\
  can you make some joke about programmer for me ?
  \tcbline
  \textbf{Assistant A} \\
  Why do programmers always mix up Christmas and Halloween?\\\\
  Because Oct 31 == Dec 25!
  \tcbline
  \textbf{Assistant B} \\
  Sure, I can make a joke about programmers. Here's a joke about programmers: \\\\
  Why did the programmer refuse to eat the two cookies he had been given at lunch? \\\\
  Because he wanted to save them for later!
  \tcbline
  \textbf{Evaluator}: \textit{Abstained} \quad \textbf{Model Preference}: \textit{None} \quad \textbf{Human Preference}: \textit{B is better} \\
  \textbf{Inter-Annotator Agreement}: $0.75$
\end{tcolorbox}